\documentclass[aspectratio=169,hyperref={bookmarks=false,urlcolor=blue,colorlinks=true},10pt]{beamer}
\usepackage{lmodern}
\usepackage[T1]{fontenc}
\usepackage{amsmath}
\usepackage{amssymb}
\usepackage{mathtools}
\usepackage{bm}
\usepackage{xcolor}
\usepackage[separate-uncertainty=true,multi-part-units=single,range-phrase=--,range-units=single]{siunitx}
\usepackage{url}
\usepackage{graphicx}
\usepackage{subcaption}
\usepackage{fnpct}
\usepackage{xspace}
\usepackage{listings}
\usepackage{algorithm2e}

\lstset{basicstyle=\ttfamily,frame=single}

\SetEndCharOfAlgoLine{}

\usecolortheme{default}

\setbeamertemplate{blocks}[rounded]
\setbeamertemplate{navigation symbols}{}

\setbeamertemplate{footline}{
  \hfill%
  \usebeamercolor[fg]{page number in head/foot}%
  \usebeamerfont{page number in head/foot}%
  \setbeamertemplate{page number in head/foot}[framenumber]%
  \usebeamertemplate*{page number in head/foot}\kern1em\vskip2pt%
}

\begin{document}

\begin{frame}
  \frametitle{Asymptotic comptuational complexity}

  \textbf{Computational complexity}: Measuring the resources required to solve a computational
  problem

  \vspace{\baselineskip}

  \pause
  \textbf{Aysmptotic}: Only care about long-term (large) behavior

  \vspace{2\baselineskip}

  \pause
  Key question: \textit{\textbf{\large If we increase the input, how much harder is the problem?}}
\end{frame}

\begin{frame}
  \frametitle{How to quantify complexity}

  Relationship between \textbf{problem size} and \textbf{resource usage} for a computational
  procedure

  \vspace{\baselineskip}

  \pause
  Usually represented as a function $f(n)$, where problem size $n$ is the input

  \vspace{\baselineskip}

  \pause
  Resources of interest are usually \textbf{time} or \textbf{space}
  \begin{itemize}
  \item Time: number of ``computation steps'' in the procedure
  \item Space: amount of state required for the procedure
  \end{itemize}

  \vspace{\baselineskip}

  \pause
  What constitutes a single ``computation step'' is axiomatic (we must decide arbitrarily, but are
  informed by reality)
\end{frame}

\begin{frame}
  \frametitle{Constant-time operations}

  \textbf{Constant-time operations}: building blocks assumed to take a single computation step

  \vspace{\baselineskip}

  \pause
  Some common examples:
  \begin{itemize}
  \item Simple arithmetic
  \item Data storage and retrieval (for fixed size)
  \item Control-flow operations (such as if-statements)
  \end{itemize}

  \vspace{\baselineskip}

  \pause
  Choice of constant-time operations requires assumptions about the model of computation

  \vspace{\baselineskip}

  \pause
  Fixed numbers of constant time operations are also constant time
\end{frame}

\begin{frame}[fragile]
  \frametitle{Big-$O$ notation: ignoring what we don't care about}

  Function: $f(n) = 5n^2 + 1000n + 10^6$

  \vspace{\baselineskip}

  \pause
  Only the \textbf{asyptotically largest} term matters
  \begin{itemize}
  \item If $n$ gets big enough, the others will be extremely small by comparison
  \item This is asymptotic analysis!
  \item $f(n) \approx 5n^2$ for very large $n$
  \end{itemize}

  \vspace{\baselineskip}

  \pause
  Ignore constant factors
  \begin{itemize}
  \item We care about the \textbf{qualitative description} of the behavior
  \item Really difficult to do bookkeeping for constant factors anyway
  \item $f(n) \approx n^2$
  \end{itemize}

  \vspace{\baselineskip}

  \pause
  Use big-$O$ notation: $f(n) = O(n^2)$
\end{frame}

\begin{frame}
  \frametitle{Example: constant vs.\ non-constant}

  \begin{columns}
    \begin{column}{0.5\textwidth}
      Problem: add two numbers $a$ and $b$

      \pause
      \vspace{\baselineskip}

      Algorithm 1:
      \begin{algorithm}[H]
        Return $a + b$
      \end{algorithm}

      \pause
      \vspace{\baselineskip}

      Algorithm 2:
      \begin{algorithm}[H]
        \For{$i$ in range(1,$b$)}{
          $a = a+1$\;
        }
        Return $a$\;
      \end{algorithm}
    \end{column}
    \begin{column}{0.5\textwidth}
      \pause
      Complexity?

      \pause
      \vspace{\baselineskip}

      Algorithm 1:
      \begin{itemize}
      \item Only one operation, regardless of the input
      \item \textbf{Constant time}: $O(1)$
      \end{itemize}

      \pause
      \vspace{\baselineskip}

      Algorithm 2:
      \begin{itemize}
      \item Number of operations proportional to $b$
      \item Therefore $b$ is the problem size
      \item Time required: $O(b)$
      \end{itemize}
    \end{column}
  \end{columns}
\end{frame}

\begin{frame}
  \frametitle{Distance between each element}

  \begin{columns}
    \begin{column}{0.5\textwidth}
      $n$ points, $x_1,x_2,\ldots,x_n$

      find $d(x_i,x_j)$ for each pair

      \pause
      \vspace{\baselineskip}

      \begin{algorithm}[H]
        \For{$i$ in range(1,$n$)}{
          \For{$j$ in range(1,$n$)}{
            $d(x_i,x_j) = \|x_i - x_j\|$\;
          }
        }
      \end{algorithm}

      \vspace{\baselineskip}

      Complexity?
      \pause
      $O(n^2)$
    \end{column}
    \begin{column}{0.5\textwidth}
      \pause
      Idea: don't recompute $d(x_j,x_i)$ after we know $d(x_i,x_j)$

      \pause
      \vspace{\baselineskip}

      \begin{algorithm}[H]
        \For{$i$ in range(1,$n$)}{
          \For{$j$ in range($i$,$n$)}{
            $d(x_i,x_j) = \|x_i - x_j\|$\;
          }
        }
      \end{algorithm}

      \vspace{\baselineskip}

      Complexity?
      \pause
      $\frac{1}{2}O(n^2) = O(n^2)$
    \end{column}
  \end{columns}
\end{frame}

\begin{frame}
  \frametitle{Does $A$ contain $x$?}

  \begin{columns}
    \begin{column}{0.5\textwidth}
      \textbf{Sorted} array $A$ with $n$ elements. Does $A$ contain $x$?

      \vspace{\baselineskip}

      \pause
      Algorithm 1:
      \begin{itemize}
      \item Check each element of $a$ in order
      \item If we find it, yes
      \item If $x>y$ for any comparison, no
      \item If we don't find it, no
      \end{itemize}

      \vspace{\baselineskip}

      \pause
      Algorithm 2:
      \begin{itemize}
      \item Get $y$, middle element of $A$
      \item If $x=y$, yes
      \item If $A$ is empty, no
      \item If $x<y$, check first half of $A$
      \item If $x>y$, check second half of $A$
      \end{itemize}
    \end{column}
    \begin{column}{0.5\textwidth}
      \pause
      Complexity (\textbf{Consider the worst case})

      \vspace{\baselineskip}

      \pause
      Algorithm 1:
      \begin{itemize}
      \item Worst case: $x$ bigger than any element
      \item Must check each element
      \item $O(n)$
      \end{itemize}

      \vspace{\baselineskip}

      \pause
      Algorithm 2:
      \begin{itemize}
      \item Worst case: $x$ not in $A$
      \item One check each iteration
      \item How many iterations? Problem size becomes $n,\frac{n}{2},\frac{n}{4},\ldots,1$
      \item $\log(n)$ steps
      \item $O(\log(n))$
      \end{itemize}
    \end{column}
  \end{columns}
\end{frame}

\begin{frame}
  \frametitle{Does $A$ contain $m$ elements?}

  \begin{columns}
    \begin{column}{0.5\textwidth}
      \textbf{Unsorted} $A$, check $x_1,x_2,\ldots,x_m$

      \pause
      \vspace{\baselineskip}

      Algorithm 3:
      \begin{itemize}
      \item For each $x_i$, use Algorithm 1
      \end{itemize}

      \pause
      \vspace{\baselineskip}

      Algorithm 4:
      \begin{itemize}
      \item Sort $A$
      \item For each $x_i$, use Algorithm 2
      \end{itemize}
    \end{column}
    \begin{column}{0.5\textwidth}
      \pause
      Complexity?

      \vspace{\baselineskip}

      \pause
      Algorithm 3:
      \begin{itemize}
      \item Do Algorithm 1 $m$ times
      \item $m O(n) = O(mn)$
      \end{itemize}

      \vspace{\baselineskip}

      \pause
      Algorithm 4:
      \begin{itemize}
      \item Sorting is generally $O(n\log(n))$
      \item Do Algorithm 2 $m$ times
      \item $m O(\log(n)) = O(m\log(n))$
      \item Total: $O((m+n)\log(n))$
      \end{itemize}
    \end{column}
  \end{columns}
\end{frame}

\begin{frame}
  \frametitle{Algorithm 3 vs.\ Algorithm 4}

  \begin{columns}[t]
    \begin{column}{0.5\textwidth}
      Suppose $m > O(n)$

      \vspace{\baselineskip}

      Then $O((m+n)\log(n)) = O(m\log(n))$

      \vspace{\baselineskip}

      Algorithm 4 better if $O(m\log(n)) < O(mn)$

      \vspace{\baselineskip}

      Always true!
    \end{column}
    \begin{column}{0.5\textwidth}
      \pause
      What about $m < O(n)$?

      \vspace{\baselineskip}

      Algorithm 3 better if $O(mn) < O(n\log(n))$

      \vspace{\baselineskip}

      Algorithm 3 better if $m < O(\log(n))$
    \end{column}
  \end{columns}

  \vspace{3\baselineskip}

  \pause
  \textbf{Conclusion:} Use Algorithm 3 only when $m < O(\log(n))$
\end{frame}

\begin{frame}[fragile]
  \frametitle{Why it matters}

  \onslide<+->

  Typical processor speed: \qty{5}{\giga\hertz}

  \begin{table}
    \begin{tabular}{|c|c|c|c|c|}
      \hline
      Complexity & Example & $n=10$ & $n=10^3$ & $n=10^6$\\
      \hline
      $O(1)$ & Add two numbers & \num{e-10} & \num{e-10} & \num{e-10}\\
      \onslide<+->{
        $O(\log(n))$ & Binary search & \num{e-10} & \num{e-9} & \num{e-9}
      }\\
      \onslide<+->{
        $O(n)$ & Add $n$ numbers & \num{e-9} & \num{e-7} & \num{e-4}
      }\\
      \onslide<+->{
        $O(n\log(n))$ & Sort $n$ numbers, FFT & \num{e-9} & \num{e-6} & \num{e-3}
      }\\
      \onslide<+->{
        $O(n^2)$ & Pairwise distance & \num{e-8} & \num{e-4} & \num{e2}
      }\\
      \onslide<+->{
        $O(n^3)$ & Matrix multiplication & \num{e-7} & \num{e-1} & \num{e8}
      }\\
      \onslide<+->{
        $O(1.1^n)$ & ``Small'' exponential & \num{e-10} & \num{e31} & \num{e41382}
      }\\
      \onslide<+->{
        $O(2^n)$ & Solve TSP & \num{e-7} & \num{e291} & \num{e301020}
      }\\
      \onslide<+->{
        $O(n!)$ & Brute-force TSP & \num{e-4} & \num{e2557} & \num{e5565699}
      }\\
      \hline
    \end{tabular}
  \end{table}

  \onslide<+->
  Age of the universe: \num{4e17} seconds
\end{frame}

\end{document}

% Local Variables:
% compile-command: "pdflatex slides.tex && pdflatex slides.tex"
% End:
